\section{Structural, sampling, and statistical diagnostics}
\subsection{Cross-arm structural features}
% AUTO-GENERATED by paper/make_revision_evidence.py.
\begin{table}[t]
\centering
\caption{Descriptive source-feature averages on common oracle-passing support. These lexical features are not physical pipeline-depth measurements.}
\label{tab:structure}
\small
\begin{tabular}{@{}lrrr@{}}
\toprule
Feature & SFT & RF & Correctness \\
\midrule
Max. multiplications/statement & \revisionclaim{StructureStatementMultSft} & \revisionclaim{StructureStatementMultRf} & \revisionclaim{StructureStatementMultCorrectness} \\
Structural accumulation score & \revisionclaim{StructureAccumulationSft} & \revisionclaim{StructureAccumulationRf} & \revisionclaim{StructureAccumulationCorrectness} \\
Structural pipeline ratio & \revisionclaim{StructurePipelineRatioSft} & \revisionclaim{StructurePipelineRatioRf} & \revisionclaim{StructurePipelineRatioCorrectness} \\
\bottomrule
\end{tabular}

\end{table}

The comparison uses \revisionclaim{StructureDesigns} designs with oracle-passing
support in all three arms. Within each design, features are averaged over
passing candidate multiplicities; designs then receive equal weight. RF has
smaller per-statement multiplier groups and larger structural accumulation and
pipeline-ratio scores than SFT. Correctness-only shifts in the opposite direction.
These source-derived features describe expressions and assignment patterns,
including patterns available to the RF reward. They do not measure routed
pipeline depth or establish an independent causal mechanism. The main paper's
implemented FIR pair connects this descriptive evidence to a physical example.

\subsection{Finite-support sampling analysis}
For a perfect-selector comparison, the empirical SFT distribution contains
every measured candidate with its multiplicity plus a zero outcome for every
failure.  If $G_d(x)$ is its empirical cumulative distribution, the maximum of
$N$ independent draws has cumulative distribution $G_d(x)^N$.  This exact
curve assumes access to every draw's measured timing-derived frequency before
selection.  It upper-bounds reranking on this finite sampled support, not on the
unknown full SFT distribution.  Draw count is not wall-clock cost: oracle calls,
uncached correct-draw implementation attempts, and content-deduplicated
implementations are reported separately.  Training cost is not amortized into
these inference-only counts.

To measure concentration, define the all-sample candidate mass
$w_{p,d,m}=c_m/n_d$ and the dominant mass $\max_m w_{p,d,m}$.  Normalizing over
correct candidates gives $\widetilde{w}_{p,d,m}=c_m/\sum_jc_j$ and effective
support
\begin{equation}
D_{p,d}=\left(\sum_m\widetilde{w}_{p,d,m}^{\two}\right)^{-\one}.
\label{eq:effective-support}
\end{equation}
Lower $D_{p,d}$ means that probability is concentrated on fewer correct strings,
which does not by itself imply better hardware. We therefore examine it
alongside measured candidate frequencies and ask whether SFT sampling already
contains an implementation at least as fast as the optimized conditional mean.
Among \studyclaim{StudyTwoSealedMechanismCommonCount} designs with
correct sampled strings under both policies, SFT's measured support contains a
candidate at least as fast as the RF conditional mean on
\studyclaim{StudyTwoSealedMechanismSupportCount}.  Mean inverse-Simpson
effective correct-string support contracts from
\studyclaim{StudyTwoSealedMechanismSftEffective} to
\studyclaim{StudyTwoSealedMechanismRfEffective}. These observations support a
shift toward faster implementations on most designs with common support,
although the finite samples and the sealed exceptions leave circuit novelty
unresolved.

On the sealed distribution, one RF draw lies between the empirical SFT
perfect-selector expectations for \revisionclaim{SelectorLowerN} draws
(\revisionclaim{SelectorLowerFmax}~MHz) and
\revisionclaim{SelectorUpperN} draws
(\revisionclaim{SelectorUpperFmax}~MHz).  At the upper integer crossing, the
SFT workload is \revisionclaim{SelectorUpperN} generations and
\revisionclaim{SelectorStreams} oracle simulation streams per design, with
\revisionclaim{SelectorRawImplementations} expected correct-draw implementation
attempts without caching or \revisionclaim{SelectorCachedImplementations}
distinct implementations with content caching.  These are empirical expected
workloads, not measured runtime or speedup.  They exclude policy-training cost
and assume perfect post-route selection on observed support; they do not imply
dominance over an unrestricted SFT search.

\subsection{Family effects, endpoint decomposition, and seed dispersion}
% AUTO-GENERATED by analyze_study2_secondary.py. DO NOT EDIT.
\begin{table*}[t]
\centering
\caption{Independent repaired-RF training-seed effects on the sealed endpoint.
Seed rows are separately reported point estimates; the seed-averaged row uses
the frozen paired, stratified 95\% bootstrap intervals.}
\label{tab:study2-rf-seeds}
\small
\begin{tabular}{lrrr}
\toprule
Estimate & Interpolation & Extrapolation & Overall \\
\midrule
RF seed 1 & \studyclaim{StudyTwoRfSeedOneInterpGain} &
\studyclaim{StudyTwoRfSeedOneExtrapGain} &
\studyclaim{StudyTwoRfSeedOneOverallGain} \\
RF seed 2 & \studyclaim{StudyTwoRfSeedTwoInterpGain} &
\studyclaim{StudyTwoRfSeedTwoExtrapGain} &
\studyclaim{StudyTwoRfSeedTwoOverallGain} \\
Seed average [95\% CI] &
\studyclaim{StudyTwoInterpRfSftGain}
[\studyclaim{StudyTwoInterpRfSftCiLow}, \studyclaim{StudyTwoInterpRfSftCiHigh}] &
\studyclaim{StudyTwoExtrapRfSftGain}
[\studyclaim{StudyTwoExtrapRfSftCiLow}, \studyclaim{StudyTwoExtrapRfSftCiHigh}] &
\studyclaim{StudyTwoOverallRfSftGain}
[\studyclaim{StudyTwoOverallRfSftCiLow}, \studyclaim{StudyTwoOverallRfSftCiHigh}] \\
\bottomrule
\end{tabular}
\end{table*}


% AUTO-GENERATED by analyze_study2_secondary.py. DO NOT EDIT.
\begin{table*}[t]
\centering
\caption{Post-primary descriptive family breakdown of the sealed
failure-penalized endpoint. I/E gives interpolation/extrapolation design counts.
Rows are descriptive point estimates; the small and unequal family cells do not
support stable family-level confidence intervals.}
\label{tab:study2-family}
\small
\begin{tabular}{lrrrrr}
\toprule
Family & Designs & I/E & SFT [MHz] & RF [MHz] & RF--SFT [MHz] \\
\midrule
FIR & \studyclaim{StudyTwoFamilyFirCount} & \studyclaim{StudyTwoFamilyFirInterpCount}/\studyclaim{StudyTwoFamilyFirExtrapCount} & \studyclaim{StudyTwoFamilyFirSft} & \studyclaim{StudyTwoFamilyFirRf} & \studyclaim{StudyTwoFamilyFirGain} \\
Reverse FIR & \studyclaim{StudyTwoFamilyReverseFirCount} & \studyclaim{StudyTwoFamilyReverseFirInterpCount}/\studyclaim{StudyTwoFamilyReverseFirExtrapCount} & \studyclaim{StudyTwoFamilyReverseFirSft} & \studyclaim{StudyTwoFamilyReverseFirRf} & \studyclaim{StudyTwoFamilyReverseFirGain} \\
Polynomial & \studyclaim{StudyTwoFamilyPolynomialCount} & \studyclaim{StudyTwoFamilyPolynomialInterpCount}/\studyclaim{StudyTwoFamilyPolynomialExtrapCount} & \studyclaim{StudyTwoFamilyPolynomialSft} & \studyclaim{StudyTwoFamilyPolynomialRf} & \studyclaim{StudyTwoFamilyPolynomialGain} \\
Recursive filter & \studyclaim{StudyTwoFamilyRecursiveCount} & \studyclaim{StudyTwoFamilyRecursiveInterpCount}/\studyclaim{StudyTwoFamilyRecursiveExtrapCount} & \studyclaim{StudyTwoFamilyRecursiveSft} & \studyclaim{StudyTwoFamilyRecursiveRf} & \studyclaim{StudyTwoFamilyRecursiveGain} \\
Median filter & \studyclaim{StudyTwoFamilyMedianCount} & \studyclaim{StudyTwoFamilyMedianInterpCount}/\studyclaim{StudyTwoFamilyMedianExtrapCount} & \studyclaim{StudyTwoFamilyMedianSft} & \studyclaim{StudyTwoFamilyMedianRf} & \studyclaim{StudyTwoFamilyMedianGain} \\
\bottomrule
\end{tabular}
\end{table*}


Table~\ref{tab:study2-family} shows how the post-primary effects vary by family.  FIR, reverse FIR, polynomial, and recursive-filter
point estimates are positive, whereas median is negative and is represented
only by extrapolation designs.  These rows are descriptive because the family
cells are small and unequal.  The primary direction is nevertheless insensitive
to two post-primary checks: the unstratified paired percentile
bootstrap gives
[\studyclaim{StudyTwoSensitivityUnstratifiedCiLow},
\studyclaim{StudyTwoSensitivityUnstratifiedCiHigh}]~MHz, and the mean effect
across leave-one-family-out omissions ranges from
\studyclaim{StudyTwoSensitivityLeaveFamilyOutMin} to
\studyclaim{StudyTwoSensitivityLeaveFamilyOutMax}~MHz.  Both checks are post-
primary sensitivities, not replacements for the frozen interval.
% AUTO-GENERATED by paper/make_revision_evidence.py.
\begin{table}[t]
\centering
\caption{Post-primary decomposition on the complete sealed draw set. Conditional frequency pools correct draws, including zero-valued implementation failures.}
\label{tab:decomposition}
\small
\begin{tabular}{@{}lrrr@{}}
\toprule
Policy & $q$ (\%) & $\mu$ (MHz) & $q\mu$ (MHz) \\
\midrule
SFT & \revisionclaim{SftCorrectness} & \revisionclaim{SftConditional} & \revisionclaim{SftPenalized} \\
RF & \revisionclaim{RfCorrectness} & \revisionclaim{RfConditional} & \revisionclaim{RfPenalized} \\
MLP & \revisionclaim{MlpCorrectness} & \revisionclaim{MlpConditional} & \revisionclaim{MlpPenalized} \\
Correctness only & \revisionclaim{CorrectnessCorrectness} & \revisionclaim{CorrectnessConditional} & \revisionclaim{CorrectnessPenalized} \\
\bottomrule
\end{tabular}

\end{table}


Table~\ref{tab:decomposition} separates functional yield from the physical
quality of oracle-passing draws.  With equal budgets across designs, pooled
$q\mu$ reproduces the primary mean exactly; this pooled $\mu$ is not an
unweighted mean of per-design conditional frequencies.  RF increases pooled
conditional frequency from \revisionclaim{SftConditional} to
\revisionclaim{RfConditional}~MHz while correctness decreases.  In contrast,
the correctness-only policy has a conditional frequency of
\revisionclaim{CorrectnessConditional}~MHz.  The observed RF gain is thus a
shift in the physical quality of passing draws, not an increase in their count.
All \revisionclaim{SeedCandidateCount} additional-seed candidate records and
\revisionclaim{SeedTrialCount} trial records pass the artifact audit and exact
offline replay of the frozen search. The RF seed mean is
\revisionclaim{SeedRfMean}~MHz, with a paired gain of
\revisionclaim{SeedRfGain}~MHz
[\revisionclaim{SeedRfLow}, \revisionclaim{SeedRfHigh}]; its descriptive
between-seed sample standard deviation is \revisionclaim{SeedRfSd}~MHz.
The correctness-only mean is \revisionclaim{SeedCorrMean}~MHz, with difference
\revisionclaim{SeedCorrGain}~MHz
[\revisionclaim{SeedCorrLow}, \revisionclaim{SeedCorrHigh}] and seed standard
deviation \revisionclaim{SeedCorrSd}~MHz.  Unstratified intervals and every
leave-one-family-out mean retain these directions.  This strengthens the
checkpoint replication evidence, but is neither a universal policy ranking nor
an independent full-pipeline repeat.  These newer-tool setup-period endpoints
are not pooled with the original timing-derived results.
